\documentclass[12pt,a4paper]{article}
\usepackage[T2A]{fontenc}
\usepackage[utf8]{inputenc}
\usepackage[russian]{babel}
\usepackage{geometry}
\geometry{left=23mm,right=20mm,top=21mm,bottom=23mm}
\usepackage{amsmath,amssymb,mathtools}
\usepackage{array,booktabs,longtable}
\usepackage{xcolor,hyperref,tikz,pdflscape}
\hypersetup{colorlinks=true,linkcolor=blue!55!black,urlcolor=blue!55!black}
\setlength{\parindent}{0pt}
\setlength{\parskip}{6pt}
\setlength{\emergencystretch}{3em}
\newcommand{\err}{\textcolor{red}{Ошибка:}}
\newcommand{\idea}{\textcolor{blue!60!black}{Ключевая идея:}}
\newcommand{\result}{\textcolor{teal!60!black}{Итог:}}
\newcommand{\mathans}[1]{\ensuremath{\displaystyle #1}}
\begin{document}
\begin{titlepage}
\thispagestyle{empty}
\vspace*{58mm}
\begin{center}
{\Huge\bfseries Ревью домашней работы\par}
\vspace{25mm}
{\large Автор: Лёвин Артём Александрович\par}
\vspace{7mm}
{\large 16 октября 2026 года\par}
\end{center}
\vfill
\begin{center}
\begin{tikzpicture}[scale=1]
\draw[blue!35!black,line width=.75pt] (0,0) .. controls (2.5,.62) and (5,-.62) .. (7.6,0);
\end{tikzpicture}
\end{center}
\end{titlepage}

\newpage
\begin{landscape}
\section*{Сводная таблица ошибок}
\small
\setlength{\tabcolsep}{3pt}
\renewcommand{\arraystretch}{1.4}
\begin{longtable}{@{}p{.055\linewidth}p{.15\linewidth}p{.26\linewidth}p{.22\linewidth}p{.235\linewidth}@{}}
\toprule
\textbf{№} & \textbf{Тема} & \textbf{Суть ошибки} & \textbf{Ваш ответ} & \textbf{Правильный ответ}\\
\midrule
\endhead
\hyperref[task:2]{2} & Тригонометрическое уравнение & После деления записан синус вместо тангенса; ответ не вытекает из этого шага. & $\pi+2\pi n$; $\frac\pi2+2\pi n$ & Те же серии, но с верным обоснованием.\\
\hyperref[task:5]{5} & Уравнение с $\sin2x$, $\cos2x$ & Неверно сложены 24 и 20, первый корень испорчен. & $\arctan 6+\pi n$; $\arctan\frac12+\pi n$ & $\arctan\frac{11}{2}+\pi n$; $\arctan\frac12+\pi n$\\
\hyperref[task:6]{6} & Логарифмы & Знак у слагаемого $\frac12$ при сложении изменён. & $-\frac72$ & $-\frac52$\\
\hyperref[task:8]{8} & Логарифм и тригонометрия & Некорректное деление на косинус, потерян корень. & $\arctan2+\pi n$ & $\frac\pi2+2\pi n$; $\pi+2\pi n$\\
\hyperref[task:9]{9} & Логарифм и тригонометрия & Неверный переход к тангенсу, один корень потерян, другой лишний. & $\pi n$ & $2\pi n$; $\frac\pi2+2\pi n$\\
\hyperref[task:10]{10} & Четвёртые степени & Рассмотрены только частные разбиения суммы; отсутствие иных корней не доказано. & $\frac\pi2+2\pi n$ & Тот же ответ, с полным доказательством.\\
\bottomrule
\end{longtable}
\normalsize
\textit{Во всех тригонометрических ответах $n\in\mathbb Z$.}
\end{landscape}

\newpage
\thispagestyle{empty}
\vspace*{.37\textheight}
\begin{center}{\Huge\bfseries Разбор ошибок}\end{center}
\newpage

\phantomsection\label{task:2}
\section*{Задание № 2}
\textbf{Текст задания.} Решите уравнение
\[\sin x-\cos x=1.\]
\textbf{Ваше решение.} При помощи формул половинного аргумента Вы верно получили
\[
2\cos\frac x2\left(\sin\frac x2-\cos\frac x2\right)=0.
\]
Далее рассмотрены два случая: $\cos\frac x2=0$ и $\sin\frac x2-\cos\frac x2=0$.
Во втором случае после указания деления на $\cos\frac x2$ записано $\sin\frac x2=1$.

\textbf{Ваш ответ.}
\[x=\pi+2\pi n\quad\text{или}\quad x=\frac\pi2+2\pi n,\quad n\in\mathbb Z.\]

\textbf{\err} Неверно записано уравнение после деления на косинус половинного аргумента. Обе серии ответа при этом верны, но вторая не обоснована приведённым шагом.

\textbf{Где именно возникает ошибка.} Последний верный шаг:
\[\sin\frac x2-\cos\frac x2=0,\qquad \cos\frac x2\ne0.\]
Вместо написанного $\sin\frac x2=1$ должно получиться
\[\frac{\sin\frac x2}{\cos\frac x2}=1,\qquad \tan\frac x2=1.\]

\textbf{Почему это неверно.} При делении равенства на ненулевое число делится \emph{каждое} слагаемое. Поэтому $\sin\frac x2$ превращается в отношение синуса к косинусу, а не остаётся синусом. Более того, из $\sin\frac x2=1$ следовали бы значения $x=\pi+4\pi n$, а не указанная Вами вторая серия.

\textbf{\idea} Сначала приравняйте каждый множитель произведения к нулю. Деление на тригонометрическую функцию допускается лишь в той ветви, где она заведомо не равна нулю.

\textbf{Исправление.} Оставляем верное разложение на множители, а во второй ветви получаем $\tan\frac x2=1$.

\textbf{Полное правильное решение.} Используем
\[\sin x=2\sin\frac x2\cos\frac x2,\quad
\cos x=\cos^2\frac x2-\sin^2\frac x2,\quad
1=\sin^2\frac x2+\cos^2\frac x2.\]
Следовательно,
\begin{align*}
2\sin\frac x2\cos\frac x2-\cos^2\frac x2+\sin^2\frac x2
 &=\sin^2\frac x2+\cos^2\frac x2,\\
2\cos\frac x2\left(\sin\frac x2-\cos\frac x2\right)&=0.
\end{align*}
Первая ветвь: $\cos\frac x2=0$, значит $\frac x2=\frac\pi2+\pi n$ и $x=\pi+2\pi n$.
Вторая ветвь: $\sin\frac x2=\cos\frac x2$. Здесь косинус не может равняться нулю, иначе синус тоже был бы нулём, что невозможно. Поэтому $\tan\frac x2=1$, откуда $\frac x2=\frac\pi4+\pi n$ и $x=\frac\pi2+2\pi n$.

\textbf{Проверка.} При $x=\pi+2\pi n$ получаем $0-(-1)=1$. При $x=\frac\pi2+2\pi n$ имеем $1-0=1$. Обе ветви получены из эквивалентного разложения.

\textbf{\result} Ответ верен, однако при делении нужно писать тангенс, а не синус.

\textbf{Задача на закрепление.} Решите уравнение $\sin x+\cos x=-1$.

\newpage
\phantomsection\label{task:5}
\section*{Задание № 5}
\textbf{Текст задания.} Решите уравнение
\[24\sin2x-7\cos2x=15.\]
\textbf{Ваше решение.} Вы правильно получили квадратное уравнение
\[4\tan^2x-24\tan x+11=0\]
и дискриминант $D=400$. Затем записано
\[t_1=\frac{24+20}{8}=6,\qquad t_2=\frac{24-20}{8}=\frac12.\]

\textbf{Ваш ответ.}
\[x=\arctan6+\pi n\quad\text{или}\quad x=\arctan\frac12+\pi n,\quad n\in\mathbb Z.\]

\textbf{\err} Арифметическая ошибка при вычислении первого корня квадратного уравнения.

\textbf{Где именно возникает ошибка.} До выражения $t_1=\frac{24+20}{8}$ вычисления верны. Первый ошибочный переход: $\frac{24+20}{8}=6$.

\textbf{Почему это неверно.} Сумма числителя равна сорока четырём, а не сорока восьми:
\[24+20=44,\qquad \frac{44}{8}=\frac{11}{2}.\]
Из-за этой арифметической ошибки вместо $\tan x=\frac{11}{2}$ получилось $\tan x=6$, что даёт постороннюю серию значений $x$.

\textbf{\idea} При замене $t=\tan x$ нужно отдельно проверить исключённый случай $\cos x=0$, а после решения квадратного уравнения ещё раз пересчитать каждую дробь.

\textbf{Исправление.} Сохраняем правильно полученное квадратное уравнение и заменяем только неверно посчитанный корень:
\[t_1=\frac{24+20}{8}=\frac{44}{8}=\frac{11}{2}.\]

\textbf{Полное правильное решение.} Если $\cos x=0$, то $\sin2x=0$, $\cos2x=-1$, и исходная левая часть равна $7$, а не $15$. Значит, этот случай корней не даёт.
При $\cos x\ne0$ положим $t=\tan x$. Тогда
\[\sin2x=\frac{2t}{1+t^2},\qquad \cos2x=\frac{1-t^2}{1+t^2}.\]
Подстановка приводит к
\begin{align*}
24\cdot\frac{2t}{1+t^2}-7\cdot\frac{1-t^2}{1+t^2}&=15,\\
48t-7+7t^2&=15+15t^2,\\
4t^2-24t+11&=0.
\end{align*}
Находим дискриминант и корни:
\[D=(-24)^2-4\cdot4\cdot11=400,\quad
 t=\frac{24\pm20}{8},\quad
 t_1=\frac{11}{2},\quad t_2=\frac12.\]
Уравнение $\tan x=a$ имеет решения $x=\arctan a+\pi n$.

\textbf{Проверка.} При $t=\frac{11}{2}$ получаем $\sin2x=\frac{44}{125}$, $\cos2x=-\frac{117}{125}$, поэтому
\[24\cdot\frac{44}{125}-7\left(-\frac{117}{125}\right)=15.\]
При $t=\frac12$ имеем $\sin2x=\frac45$, $\cos2x=\frac35$, поэтому $24\cdot\frac45-7\cdot\frac35=15$.

\textbf{\result} Правильные серии:
\[\boxed{x=\arctan\frac{11}{2}+\pi n\quad\text{или}\quad x=\arctan\frac12+\pi n,\ n\in\mathbb Z.}\]

\textbf{Задача на закрепление.} Решите уравнение $10\sin2x-6\cos2x=10$.

\newpage
\phantomsection\label{task:6}
\section*{Задание № 6}
\textbf{Текст задания.} Известно, что $\log_b a=-2$, где $a,b>0$, $a,b\ne1$. Вычислите
\[\log_a\left(b^6\sqrt a\right).\]

\textbf{Ваше решение.} Вы правильно применили свойства логарифмов:
\[\log_a(b^6\sqrt a)=6\log_a b+\frac12,\quad
\log_a b=\frac1{\log_b a}=-\frac12.\]
Затем из $-3+\frac12$ был получен результат $-\frac72$.

\textbf{Ваш ответ.} $-\frac72$ (минус три с половиной).

\textbf{\err} При приведении к общему знаменателю изменился знак положительного слагаемого.

\textbf{Где именно возникает ошибка.} Последняя правильная запись:
\[6\left(-\frac12\right)+\frac12=-3+\frac12.\]
Неверный итог: $-\frac72$.

\textbf{Почему это неверно.} При сложении отрицательного и положительного чисел модули вычитаются, а знак остаётся от числа с большим модулем. Здесь
\[-3+\frac12=-\frac62+\frac12=\frac{-6+1}{2}=-\frac52.\]
Число $-\frac72$ получилось бы из $-3-\frac12$, но в исходном выражении стоит знак «плюс».

\textbf{\idea} Сначала применяйте свойства логарифмов, затем приводите обычные дроби к одному знаменателю без изменения знаков.

\textbf{Исправление.} От верной записи $-3+\frac12$ приходим к $-\frac52$.

\textbf{Полное правильное решение.} Пользуемся тождествами $\log_a(uv)=\log_a u+\log_a v$, $\log_a(u^k)=k\log_a u$ и $\log_a b=1/\log_b a$:
\begin{align*}
\log_a\left(b^6\sqrt a\right)
 &=\log_a b^6+\log_a a^{1/2}\\
 &=6\log_a b+\frac12\log_a a\\
 &=6\cdot\frac1{-2}+\frac12\\
 &=-3+\frac12=-\frac52.
\end{align*}
Условия $a,b>0$ и $a,b\ne1$ обеспечивают существование использованных логарифмов; $b^6\sqrt a>0$.

\textbf{Проверка.} $-\frac52=-2{,}5$, а $-3+0{,}5=-2{,}5$.

\textbf{\result} $\boxed{-\frac52}$.

\textbf{Задача на закрепление.} Известно, что $\log_c d=-3$, где $c,d>0$, $c,d\ne1$. Вычислите $\log_d(c^4\sqrt d)$.

\newpage
\phantomsection\label{task:8}
\section*{Задание № 8}
\textbf{Текст задания.} Решите уравнение
\[\log_2\left(\sqrt2\cos\left(x-\frac\pi2\right)-\cos\left(x-\frac\pi4\right)\right)=-\frac12.\]

\textbf{Ваше решение.} Вы верно раскрыли формулы приведения и сложения углов и дошли до
\[\sin x-\cos x=1.\]
После записи «делим на $\cos x\ne0$» появилось $\tan x-1=1$, затем $\tan x=2$.

\textbf{Ваш ответ.} $x=\arctan2+\pi n$, $n\in\mathbb Z$.

\textbf{\err} При делении уравнения на косинус не разделена правая часть; дополнительно исключены допустимые значения с нулевым косинусом.

\textbf{Где именно возникает ошибка.} Последний верный шаг: $\sin x-\cos x=1$. Неверный переход: $\tan x-1=1$.

\textbf{Почему это неверно.} При $\cos x\ne0$ деление даёт
\[\frac{\sin x}{\cos x}-\frac{\cos x}{\cos x}=\frac1{\cos x},\qquad
\tan x-1=\frac1{\cos x},\]
а не $\tan x-1=1$. Кроме того, при $\cos x=0$ делить вообще нельзя. Именно среди таких значений есть решение $x=\frac\pi2+2\pi n$.

\textbf{\idea} После перехода от логарифма к его аргументу решайте полученное тригонометрическое уравнение способом, не требующим деления на переменную функцию.

\textbf{Исправление.} Оставляем верное $\sin x-\cos x=1$ и переписываем его как $\sqrt2\sin\left(x-\frac\pi4\right)=1$.

\textbf{Полное правильное решение.} Для существования логарифма его аргумент должен быть положителен. Так как основание равно двум,
\[\log_2 A=-\frac12\quad\Longleftrightarrow\quad A=2^{-1/2}=\frac1{\sqrt2}>0.\]
Из формул сложения:
\[\cos\left(x-\frac\pi2\right)=\sin x,\quad
\cos\left(x-\frac\pi4\right)=\frac{\sqrt2}{2}(\cos x+\sin x).\]
Следовательно,
\begin{align*}
\sqrt2\sin x-\frac{\sqrt2}{2}(\cos x+\sin x)&=\frac{\sqrt2}{2},\\
\sin x-\cos x&=1,\\
\sin\left(x-\frac\pi4\right)&=\frac{\sqrt2}{2}.
\end{align*}
Последнее уравнение имеет две серии:
\[x-\frac\pi4=\frac\pi4+2\pi n\quad\text{или}\quad
x-\frac\pi4=\frac{3\pi}{4}+2\pi n.\]
Поэтому $x=\frac\pi2+2\pi n$ или $x=\pi+2\pi n$.

\textbf{Проверка.} Для первой серии $\sin x-\cos x=1-0=1$, для второй $0-(-1)=1$. Следовательно, в обоих случаях исходный аргумент логарифма равен $1/\sqrt2>0$, а логарифм равен $-\frac12$.

\textbf{\result} $\boxed{x=\frac\pi2+2\pi n\ \text{или}\ x=\pi+2\pi n,\ n\in\mathbb Z.}$

\textbf{Задача на закрепление.} Решите уравнение $\log_2\left(\sqrt2\cos(x-\frac\pi2)-\cos(x-\frac\pi4)\right)=0$.

\newpage
\phantomsection\label{task:9}
\section*{Задание № 9}
\textbf{Текст задания.} Решите уравнение
\[\log_2\left(\sqrt2\cos\left(x-\frac\pi2\right)+\cos\left(x+\frac\pi4\right)\right)=-\frac12.\]

\textbf{Ваше решение.} Верно преобразовав аргумент логарифма, Вы получили
\[\sin x+\cos x=1.\]
Далее при делении на $\cos x\ne0$ было записано $\tan x=0$.

\textbf{Ваш ответ.} $x=\pi n$, $n\in\mathbb Z$.

\textbf{\err} Ошибочный переход к тангенсу привёл к лишним значениям и потере правильной серии.

\textbf{Где именно возникает ошибка.} Последнее верное равенство: $\sin x+\cos x=1$. При делении на ненулевой косинус должно быть
\[\tan x+1=\frac1{\cos x},\]
а не $\tan x=0$.

\textbf{Почему это неверно.} Правая единица при делении на $\cos x$ превращается в $1/\cos x$, а не остаётся неизменной. Кроме того, $x=\pi+2\pi n$ из Вашего ответа неверен: для него $\sin x+\cos x=-1$, и аргумент исходного логарифма отрицателен. Значения $x=\frac\pi2+2\pi n$, напротив, дают верное равенство, но были потеряны при делении на косинус.

\textbf{\idea} Не нужно делить на косинус: сумму синуса и косинуса можно преобразовать в один синус со сдвигом.

\textbf{Исправление.} Возвращаемся к уравнению $\sin x+\cos x=1$ и используем $\sqrt2\sin(x+\frac\pi4)=1$.

\textbf{Полное правильное решение.} Пусть $A$ --- выражение под логарифмом. Уравнение логарифмов равносильно $A=1/\sqrt2$, что автоматически обеспечивает $A>0$.
\[\cos\left(x-\frac\pi2\right)=\sin x,\qquad
\cos\left(x+\frac\pi4\right)=\frac{\sqrt2}{2}(\cos x-\sin x).\]
Тогда
\begin{align*}
\sqrt2\sin x+\frac{\sqrt2}{2}(\cos x-\sin x)&=\frac{\sqrt2}{2},\\
\sin x+\cos x&=1,\\
\sin\left(x+\frac\pi4\right)&=\frac{\sqrt2}{2}.
\end{align*}
Отсюда
\[x+\frac\pi4=\frac\pi4+2\pi n\quad\text{или}\quad
x+\frac\pi4=\frac{3\pi}{4}+2\pi n.\]
Поэтому $x=2\pi n$ или $x=\frac\pi2+2\pi n$.

\textbf{Проверка.} При $x=2\pi n$ имеем $\sin x+\cos x=0+1=1$; при $x=\frac\pi2+2\pi n$ получаем $1+0=1$. В обоих случаях аргумент логарифма равен $1/\sqrt2>0$.

\textbf{\result} $\boxed{x=2\pi n\ \text{или}\ x=\frac\pi2+2\pi n,\ n\in\mathbb Z.}$

\textbf{Задача на закрепление.} Решите уравнение $\log_2\left(\sqrt2\cos(x-\frac\pi2)+\cos(x+\frac\pi4)\right)=0$.

\newpage
\phantomsection\label{task:10}
\section*{Задание № 10}
\textbf{Текст задания.} Решите уравнение
\[\sin^4x+(\sin x-2)^4=2\]
и обоснуйте, почему других корней нет.

\textbf{Ваше решение.} Вы рассмотрели частный вариант, когда оба слагаемых равны единице, а также вариант, когда первое равно нулю, а второе равно двум. В первом варианте получено $\sin x=1$ и верная серия значений $x$; второй вариант отвергнут.

\textbf{Ваш ответ.} $x=\frac\pi2+2\pi n$, $n\in\mathbb Z$.

\textbf{\err} Верный ответ найден, но требование доказать отсутствие других корней не выполнено.

\textbf{Где именно возникает ошибка.} Первое необоснованное допущение --- замена условия «сумма двух четвёртых степеней равна двум» рассмотрением только случаев «один плюс один» и «ноль плюс два».

\textbf{Почему это неверно.} Вообще из $A+B=2$ не следует, что обязательно $A=B=1$ или одно из чисел равно нулю. Например, положительные числа $\frac12$ и $\frac32$ тоже дают сумму два. Поэтому перечисленные Вами случаи не исчерпывают всех возможных значений слагаемых, даже если они четвёртые степени. Нужно обосновать, что в \emph{данном} уравнении иной вариант невозможен.

\textbf{\idea} Используем ограничение $-1\le\sin x\le1$ и представим обе четвёртые степени симметрично относительно числа один. Тогда появится сумма неотрицательных выражений.

\textbf{Исправление.} Положим $u=1-\sin x$. Так как $\sin x\le1$, то $u\ge0$, а
\[\sin x=1-u,\qquad \sin x-2=-1-u.\]
Получаем уравнение, в котором все добавочные к двойке слагаемые неотрицательны.

\textbf{Полное правильное решение.} Подставим выражения через $u$:
\[(1-u)^4+(-1-u)^4=2.\]
Четвёртая степень не зависит от знака основания, поэтому
\[(1-u)^4+(1+u)^4=2.\]
Раскрываем скобки по формуле четвёртой степени суммы и разности:
\begin{align*}
(1-u)^4&=1-4u+6u^2-4u^3+u^4,\\
(1+u)^4&=1+4u+6u^2+4u^3+u^4.
\end{align*}
При сложении нечётные степени сокращаются:
\[2+12u^2+2u^4=2,\qquad 2u^2(6+u^2)=0.\]
Для любого действительного $u$ верно $u^2\ge0$ и $6+u^2>0$, поэтому произведение равно нулю тогда и только тогда, когда $u=0$.
Следовательно, $1-\sin x=0$, то есть $\sin x=1$, откуда
\[x=\frac\pi2+2\pi n,\qquad n\in\mathbb Z.\]
Другие значения $u$ невозможны, что доказывает полноту ответа.

\textbf{Проверка.} При $\sin x=1$ левая часть исходного уравнения равна $1^4+(1-2)^4=1+1=2$. При $u\ne0$ величина $12u^2+2u^4$ строго положительна, поэтому исходная сумма строго больше двух.

\textbf{\result} $\boxed{x=\frac\pi2+2\pi n,\ n\in\mathbb Z.}$ Ответ совпадает с Вашим, но теперь отсутствие других корней строго доказано.

\textbf{Задача на закрепление.} Решите уравнение $(\sin x+1)^4+(\sin x-1)^4=2$ и обоснуйте, почему других корней нет.

\vfill
\begin{center}\footnotesize Лёвин Артём Александрович эксклюзивно для Николь Саркисьянц\end{center}
\end{document}